\scanpage{51}
% archive foliation: 50
nienne sur le corps des fractions \deleted{de} L de K ?

On notera que même dams le cas où $A=W(k)$, k alg. clos, le foncteur en question n'est pas \deleted{\unsure{essentiellem}} essentiellement surjectif.

Notons le

\uline{Problème} 5. \inserted{(A anneau de val. discrète complet d'inégales car.,k parf)} Expliciter le foncteur naturel qui va de la sous-catégorie pleine précédente de la catégorie des couples \deleted{d} $(M,\mathrm{Fil})$ d'un module de Dieudonné M sur \deleted{W} k (i.e. \deleted{sur} module sur W) et d'une filtration à deux crans de $M\otimes_W L$, vers les modules galoisiens sur L (\unsure{représent tions|représentations} \inserted{\unsure{vectoorielles|vectorielles} des dimension finie sur $\underline{Q}_p$ de} de $\mathrm{Gal}(\overline{L}/L)$)\deleted{dans des \unsure{repré} modules libres de type fini} Peut-on même définir ce foncteur \deleted{sur la sous-catégorie} de façon naturelle sans se borner à cette sous-catégorie pleine,en imposant seulement des conditions de rangs sur le gradué, mettons ?

\uline{Problème} 6. Revenant à un A comme dans Pb 4, peut-on trouver un foncteur \deleted{d'une} de la catégorie des couples $(M,\mathrm{Fil}^{1})$ d'un isocristal de Dieudonné \inserted{admissible} sur $A_{\unsure{o|0}}=A/pA$ (ou encore, sur $A_{\unsure{o|0}\,\mathrm{red}}\ \deleted{\unsure{\mathrm{perf}}}$) et d'une filtration en deux crans par modules localement libres de $(M_A)_p$, ayant le rang qu'il faut\deleted{\unsure{,}} (et soumis éventuellement à d'autres conditions ouvertes,\inserted{au sens Zar} assurant que le couple envisagé provient d'un isogroupe de BT sur A),vers les faisceaux p-adiques constants tordus sur $U=\mathrm{Spec}(A_p)$, qui exprime le foncteur du théorème (VII 10.8.2) précédent.

\deleted{4) Question des relations entre cohomologie cristalline et p-adique.}

\deleted{\deleted{Problème} \uline{Conjecture} 7. On peut, pour tout anneau de valuation discrète complet V d'inégales caractéristiques, de corps \deleted{des fractions} \unsure{rismduel|résiduel} parfait (pour simplifier), de corps des fractions L, trouver un foncteur canoniques sur une sous-catégorie pleine convenable des couples $(M,\mathrm{Fil}^{\deleted{\unsure{1}}})$ \deleted{\unsure{envisagé|envisagés}} dans le Pb 5 vers les}

